\section{The density matrix}

\SourcePage{58}{61}
Describing a system by its wave function represents the fullest description
that quantum mechanics allows, in the sense indicated at the end of \secref{1}.

States not admitting such a description arise when the system considered is
part of a larger closed system. Suppose the closed system as a whole is in a
state described by $\Psi(q,x)$, where $x$ denotes the coordinates of the
system considered and $q$ the remaining coordinates of the closed system. In
general this function does not factor into functions of $x$ alone and $q$
alone, so that
\SourcePage{59}{62}
the system has no wave function of its own\footnote{For $\Psi(q,x)$ to factor
in this way at a given instant, the measurement that produced the state must
describe completely and separately the system considered and the remainder of
the closed system. For the factorization to persist, these two parts must also
not interact (see \secref{2}). Neither condition is assumed here.}.

Let $f$ be a physical quantity pertaining to our system. Its operator acts on
the $x$ coordinates only, not on $q$. Its mean value is
\begin{equation}
 \bar f=\iint\Psi^*(q,x)\hat f\Psi(q,x)\,dq\,dx.
 \tag{14.1}
\end{equation}
Introduce the function
\begin{equation}
 \rho(x,x')=\int\Psi(q,x)\Psi^*(q,x')\,dq,
 \tag{14.2}
\end{equation}
where integration is only over $q$; it is called the system's \textit{density
matrix}. From (14.2),
\begin{equation}
 \rho^*(x,x')=\rho(x',x).
 \tag{14.3}
\end{equation}
The ``diagonal elements''
\[
 \rho(x,x)=\int|\Psi(q,x)|^2\,dq
\]
determine the coordinate probability distribution.

In terms of the density matrix,
\begin{equation}
 \bar f=\int[\hat f\rho(x,x')]_{x'=x}\,dx.
 \tag{14.4}
\end{equation}
Here $\hat f$ acts only on the variables $x$ in $\rho(x,x')$; after its action
has been evaluated, one sets $x'=x$. Thus the density matrix determines the
mean of every quantity characterizing the system and hence also the
probabilities of its possible values. A state having no wave function can
therefore be described by a density matrix. The density matrix contains none
of the coordinates $q$ that do not belong
\SourcePage{60}{63}
to the system, although it naturally depends in substance on the state of the
closed system as a whole.

The density-matrix description is the most general form of
quantum-mechanical description of systems. The wave-function description is the special
case $\rho(x,x')=\Psi(x)\Psi^*(x')$. There is an important distinction. For a
state possessing a wave function (a \textit{pure} state), there is always a
complete system of measurements whose results can be predicted with certainty
(mathematically, $\Psi$ is an eigenfunction of some operator). For states
possessing only a density matrix (\textit{mixed} states), no complete system
of measurements has uniquely predictable results.

Suppose the system is closed, or became closed at some instant. We derive an
equation governing the time dependence of its density matrix, analogous to the
wave equation. The desired linear differential equation for $\rho(x,x',t)$
must also hold in the special case
\[
 \rho(x,x',t)=\Psi(x,t)\Psi^*(x',t).
\]
Differentiating and using (8.1),
\[
 \begin{aligned}
 i\hbar\frac{\partial\rho}{\partial t}
 &=i\hbar\Psi^*(x',t)\frac{\partial\Psi(x,t)}{\partial t}
 +i\hbar\Psi(x,t)\frac{\partial\Psi^*(x',t)}{\partial t}\\
 &=\Psi^*(x',t)\hat H\Psi(x,t)
 -\Psi(x,t)\hat H'^*\Psi^*(x',t).
 \end{aligned}
\]
Here $\hat H$ acts on functions of $x$, while $\hat H'$ is the same operator
acting on functions of $x'$. Bringing the other factors under the respective
operator signs gives
\begin{equation}
 i\hbar\frac{\partial\rho(x,x',t)}{\partial t}
 =(\hat H-\hat H'^*)\rho(x,x',t).
 \tag{14.5}
\end{equation}
Let $\Psi_n(x,t)$ be the stationary-state wave functions. Expanding the density
matrix in these functions gives
\SourcePage{61}{64}
the double series
\begin{equation}
 \rho(x,x',t)=\sum_m\sum_n a_{mn}\psi_n^*(x')\psi_m(x)
 \exp\left[\frac{i}{\hbar}(E_n-E_m)t\right].
 \tag{14.6}
\end{equation}
This plays the same role for the density matrix as (10.3) does for the wave
function. In place of the single set $a_n$ there is the double set $a_{mn}$,
which, like the density matrix, is Hermitian:
\begin{equation}
 a_{mn}=a_{nm}^*.
 \tag{14.7}
\end{equation}
Substitution in (14.4) gives
\[
 \bar f=\sum_m\sum_n a_{mn}\int\Psi_n^*(x,t)\hat f\Psi_m(x,t)\,dx,
\]
or
\begin{equation}
 \bar f=\sum_m\sum_n a_{mn}f_{nm}(t)
 =\sum_m\sum_n a_{mn}f_{nm}
 \exp\left[\frac{i}{\hbar}(E_n-E_m)t\right].
 \tag{14.8}
\end{equation}
This is analogous to (11.1)\footnote{The quantities $a_{mn}$ constitute the
density matrix in the energy representation. Description of states by such a
matrix was introduced independently by Landau and Bloch (\textit{F. Bloch}) in
1927.}.

The $a_{mn}$ must satisfy inequalities. Since the diagonal elements
$\rho(x,x)$ are probabilities, they must be positive. Hence the quadratic form
\[
 \sum_n\sum_m a_{mn}\xi_m\xi_n^*
\]
must be positive definite for arbitrary complex $\xi_n$. The standard
conditions for quadratic forms therefore apply; in particular,
\begin{equation}
 a_{nn}>0,
 \tag{14.9}
\end{equation}
\SourcePage{62}{65}
and for every $a_{mm}$, $a_{nn}$, $a_{mn}$,
\begin{equation}
 a_{mm}a_{nn}\geq|a_{mn}|^2.
 \tag{14.10}
\end{equation}
The pure case, in which the density matrix is a product of functions,
corresponds to
\begin{equation}
 a_{mn}=a_ma_n^*.
 \tag{14.11}
\end{equation}
A simple criterion distinguishes a pure from a mixed state. In the pure case,
\[
 (a^2)_{mn}=\sum_k a_{mk}a_{kn}
 =\sum_k a_ma_k^*a_ka_n^*=a_ma_n^*\sum_k|a_k|^2=a_ma_n^*,
\]
and therefore
\begin{equation}
 (a^2)_{mn}=a_{mn}.
 \tag{14.12}
\end{equation}
Thus the square of the density matrix coincides with the matrix itself.
